\documentclass{beamer}
\usepackage[utf8]{inputenc}
\usepackage[T1]{fontenc}
\usepackage[english]{babel}
\usepackage{hyperref}
\usepackage{outlines}
\usepackage{graphicx}
\usepackage{algorithm}
\usepackage{algorithmic}
\usepackage{amsmath,amssymb}
\usepackage{moresize}
\usepackage{tikz}
\usetikzlibrary{overlay-beamer-styles}
\usepackage{microtype}
\usepackage[
  backend=biber,
  style=alphabetic,
]{biblatex}
\usepackage{latexsym,multicol,booktabs,calligra}
\usepackage{listings,stackengine}
\usepackage[table,dvipsnames]{xcolor}
\usepackage{cancel}
\renewcommand{\figurename}{Figure}
\renewcommand{\algorithmname}{Algorithm}

\usepackage{SHU}
\setbeamersize{text margin left=7mm, text margin right=7mm}
\definecolor{deepblue}{rgb}{0,0,0.5}
\definecolor{deepred}{rgb}{0.6,0,0}
\definecolor{deepgreen}{rgb}{0,0.5,0}
\definecolor{halfgray}{gray}{0.55}
\definecolor{nus-orange}{RGB}{239,124,0}
\definecolor{nus-blue}{RGB}{0,61,124}

\lstset{
  language=SQL,
  basicstyle=\ttfamily\tiny,
  keywordstyle=\bfseries\color{nus-blue},
  emphstyle=\ttfamily\color{nus-blue},
  stringstyle=\color{deepgreen},
  numbers=left,
  numberstyle=\small\color{halfgray},
  rulesepcolor=\color{nus-orange},
  frame=shadowbox,
  showstringspaces=false,
  breaklines=true,
  breakatwhitespace=true,
  keepspaces=true,
  columns=flexible,
  upquote=true
}

\author{\href{https://pratik2358.github.io/}{Pratik Karmakar}}
\title{IT5008: Tutorial 6 — Stored Procedures and Triggers}
\institute{
  School of Computing,\\
  National University of Singapore
}
\date{AY25/26 S1}

\begin{document}

\begin{frame}
  \titlepage
  \begin{figure}[htpb]
    \begin{center}
      \includegraphics[keepaspectratio, scale=0.18]{nus-logo.png}
    \end{center}
  \end{figure}
\end{frame}

\section{Setup}
\begin{frame}{Scenario}
\small
Students at the \textbf{National University of Ngendipura (NUN)} buy, lend, and borrow books.\\
NUNStA commissions \emph{Apasaja Private Limited} to implement an online book exchange system that records:
\begin{itemize}\itemsep3pt
  \item Students: name, faculty, department, \textbf{email}, join year.
  \item Books: title, authors, publisher, edition, \textbf{ISBN10}, \textbf{ISBN13}.
  \item Loans: \texttt{borrowed} date, \texttt{returned} date (\texttt{NULL} if active).
\end{itemize}
Auditing preserves records of graduated students and copies with loans.\\
This tutorial uses the schema/data from “Creating and Populating Tables.”
\end{frame}

\section{Questions}
\begin{frame}{Questions}
\footnotesize
\begin{enumerate}
  \item Stored Functions and Procedures
  \begin{itemize}
    \item[(a)] Implement \texttt{borrow\_book} (function/procedure) to check availability of a copy, insert a loan, and return/raise a message.
    \item Scenario: Adeline Wong (\texttt{awong007@msn.com}) tries to borrow 3 copies of “Applied Calculus” (\texttt{ISBN13=978-0470170526}).
  \end{itemize}
  \item Triggers
  \begin{itemize}
    \item[(a)] Create a trigger that enforces: a student may have at most 3 active loans (local strategy).
    \item[(b)] Create a trigger that ensures globally no student exceeds 3 active loans (global strategy).
  \end{itemize}
\end{enumerate}
\end{frame}

\section{Stored Procedures}
\begin{frame}[fragile]{1(a).\ Function \texttt{borrow\_book}}
\vspace{-0.3em}
\begin{columns}[t,onlytextwidth]
  \column{0.49\textwidth}
\begin{lstlisting}[numbers=none]
CREATE OR REPLACE FUNCTION borrow_book (
  borrower_email VARCHAR(256),
  isbn13 CHAR(14),
  borrow_date DATE
) RETURNS TEXT AS $$
DECLARE
  available_copy RECORD;
BEGIN
  SELECT * INTO available_copy
  FROM copy c
  WHERE c.book = isbn13
    AND NOT EXISTS (
      SELECT 1 FROM loan l
      WHERE l.book=c.book
        AND l.copy=c.copy
        AND l.owner=c.owner
        AND l.returned IS NULL)
  LIMIT 1;
\end{lstlisting}

  \column{0.49\textwidth}
\begin{lstlisting}[numbers=none]
  IF NOT FOUND THEN
    RETURN 'No available copies of '||isbn13;
  ELSE
    INSERT INTO loan (borrower, owner, book, copy, borrowed)
    VALUES (borrower_email, available_copy.owner,
            available_copy.book, available_copy.copy, borrow_date);
    RETURN 'Book '||isbn13||' borrowed by '||borrower_email;
  END IF;
END;
$$ LANGUAGE plpgsql;
\end{lstlisting}
\end{columns}
\end{frame}

\begin{frame}[fragile]{1(a).\ Procedure \texttt{borrow\_book}}
\vspace{-0.3em}
\begin{columns}[t,onlytextwidth]
  \column{0.49\textwidth}
\begin{lstlisting}[numbers=none]
CREATE OR REPLACE PROCEDURE borrow_book (
  borrower_email VARCHAR(256),
  isbn13 CHAR(14),
  borrow_date DATE
) AS $$
DECLARE
  available_copy RECORD;
BEGIN
  SELECT * INTO available_copy
  FROM copy c
  WHERE c.book = isbn13
    AND NOT EXISTS (
      SELECT 1 FROM loan l
      WHERE l.book=c.book
        AND l.copy=c.copy
        AND l.owner=c.owner
        AND l.returned IS NULL)
  LIMIT 1;
\end{lstlisting}

  \column{0.49\textwidth}
\begin{lstlisting}[numbers=none]
  IF NOT FOUND THEN
    RAISE NOTICE 'No copies of %', isbn13;
  ELSE
    INSERT INTO loan (borrower, owner, book, copy, borrowed)
    VALUES (borrower_email, available_copy.owner,
            available_copy.book, available_copy.copy, borrow_date);
    RAISE NOTICE 'Book % borrowed by %', isbn13, borrower_email;
  END IF;
END;
$$ LANGUAGE plpgsql;
\end{lstlisting}
\end{columns}
\end{frame}

\section{Triggers}
\begin{frame}[fragile]{2(a).\ Local Loan Limit Trigger}
\textbf{Max 3 active loans per student.}
\begin{lstlisting}
CREATE OR REPLACE FUNCTION check_local_loan_limit()
RETURNS TRIGGER AS $$
DECLARE active_loan_count INT;
BEGIN
  SELECT COUNT(*) INTO active_loan_count
  FROM loan l
  WHERE l.borrower = NEW.borrower
    AND l.returned ISNULL;

  IF active_loan_count >= 3 THEN
    RETURN NULL; -- prevent insert
  ELSE
    RETURN NEW;
  END IF;
END;
$$ LANGUAGE plpgsql;

CREATE TRIGGER enforce_local_loan_limit_insert
BEFORE INSERT ON loan
FOR EACH ROW EXECUTE FUNCTION check_local_loan_limit();
\end{lstlisting}
\end{frame}

\begin{frame}[fragile]{2(b).\ Global Loan Limit Trigger}
\textbf{Ensure no student exceeds 3 active loans.}
\begin{lstlisting}
CREATE OR REPLACE FUNCTION check_global_loan_limit()
RETURNS TRIGGER AS $$
DECLARE violating_student RECORD;
BEGIN
  SELECT l.borrower INTO violating_student
  FROM loan l
  WHERE l.returned ISNULL
  GROUP BY l.borrower
  HAVING COUNT(*) > 3;

  IF violating_student IS NOT NULL THEN
    RAISE EXCEPTION '% exceeds loan limit', violating_student;
  ELSE
    RETURN NEW;
  END IF;
END;
$$ LANGUAGE plpgsql;

CREATE TRIGGER enforce_global_loan_limit
AFTER INSERT OR UPDATE ON loan
FOR EACH ROW EXECUTE FUNCTION check_global_loan_limit();
\end{lstlisting}
\end{frame}

\section{Guidelines}
\begin{frame}{Guidelines \& Remarks}
\small
\begin{itemize}\itemsep4pt
  \item Functions return a value; procedures raise notices or exceptions.
  \item Local vs Global strategies: local checks per row; global checks across all rows.
  \item Use \textbf{BEFORE} triggers with \texttt{RETURN NULL} to cancel insertions.
  \item Use \textbf{AFTER} triggers with \texttt{RAISE EXCEPTION} for global consistency.
  \item Dropping schema (\texttt{DROP SCHEMA ... CASCADE}) ensures removal of triggers/functions.
\end{itemize}
\end{frame}

\begin{frame}
\begin{center}
Questions?\\
Drop a mail at: pratik.karmakar@u.nus.edu
\end{center}
\end{frame}

\end{document}